% English translation of capitulos/03-metodologia.tex, made by Claude Code (Claude Fable 5.2) on 24 November 2026.
% Pending review by the author. When the Portuguese file changes, update this one (scripts/verificar_traducao.py).
\section{Methodology}\label{sec:metodologia}
This section describes what was carried out up to 17 November 2026 and indicates what is left for the final version. The verb
tenses follow this distinction. The word ``pilot'' has three uses in the project records, which receive
distinct names here: the six reference studies of the initial analysis, the two screening
calibration samples and the full-text measurement.

\subsection{Type of study}
This work is a declared rapid review~\citep{cartaxo2020rapid}, reported according to PRISMA
2020~\citep{page2021prisma} where applicable. The original protocol defined a systematic literature
review according to \citet{kitchenham2015ebse}, complemented by
snowballing~\citep{wohlin2014snowballing}, and provided for a fallback to a rapid review if the former
proved unfeasible. The author adopted this fallback on 3 November 2026, before the batch screening, based on the effort
projection of Section~\ref{sec:calibracao} and with the agreement of the advisor.

Cartaxo et al. allow steps to be omitted or simplified, provided that the concessions are stated in the protocol and the
report warns about the limitations. The concessions adopted here are choices of this project, and not part of the
definition of a rapid review: requiring the task terms in the title, without broad snowballing, and reading
the studies from the text extracted from the PDFs. Since 8 November 2026, the version for the examining committee also leaves out the studies
that are only on arXiv, with no indication of peer review (Section~\ref{sec:textocompleto}). The execution is largely done by AI, with checking by the
author at defined points, and every inclusion is confirmed by the author.

\subsection{Research questions and criteria}
The research questions are: RQ1, which prompting techniques are evaluated and what evidence of
effectiveness is presented, under which models, tasks, metrics and conditions; RQ2, which technical
risks are identified and which mitigation strategies are proposed or evaluated; and RQ3, how to
relate techniques, tasks and conditions in a useful taxonomy.

The selection criteria are in Table~\ref{tab:criterios}. The protocol used only open-access texts. On 5 November 2026, after the title and abstract screening,
the author started to accept also texts obtained through the CAPES journals portal (Portal de Periódicos da CAPES), with UFMG access,
to reduce availability bias. The change is reported as a deviation; the remaining bias is a threat to validity.

The application of I2 at the full-text stage was defined on 3 November 2026, after the measurement of
Section~\ref{sec:calibracao}, and not from the beginning of the protocol. The measurement brought up cases that the draft
of the prompt (tc-v1) did not resolve: an observational study that analyses prompts without measuring their effectiveness
(S2261), a workflow that uses prompting techniques without isolating their effect (S306) and a decoding
strategy with no prompting technique evaluated (S2128), for which the author used X3, and the models, X4.
Under the rule decided, recorded in version tc-v2, the study must report a code result obtained by
prompting, with the prompt strategy described, without requiring a prompting comparator. Code as an
intermediate step is eligible if it is evaluated. An observational study that only describes the use, with no result
about the prompts, goes to the author as a doubt. An evaluated intervention that is not prompting is excluded as X4,
which broadens this code at the full-text stage; at the title and abstract stage, tr-v2 keeps the previous wording, with which
it was calibrated. The change is reported as a deviation from the protocol.

With tc-v2, the two models reread in new contexts, from the extracted text, S2261, S306, S2128 and S427,
the latter included to check whether the broadening of X4 would rule out a study that compares prompting with fine-tuning.
The two answered doubt, include, exclude by X4 and include, respectively. Since the rule was written
with these cases known, the agreement shows that it is applicable to them, but does not validate it in other studies.

\begin{table}[!htb]
\caption{Selection criteria}\label{tab:criterios}
\centering\footnotesize
\begin{tabular}{@{}lp{0.82\columnwidth}@{}}
\hline
I1 & Primary study.\\
I2 & Applies prompting of a pre-trained model to generation, translation, refactoring or debugging/repair; tests are included when they support these tasks.\\
I3 & Empirical evaluation or sufficiently described case study.\\
I4 & First public version between 1 January 2021 and 7 September 2026.\\
I5 & Text in Portuguese or English.\\
I6 & Full text in open or institutional access (CAPES/UFMG).\\
\hline
X1 & Duplicate or overlapping version of a study already recorded.\\
X2 & No technical detail to identify technique, task and evaluation.\\
X3 & Task out of scope (e.g., only summarization, code review or questions about code).\\
X4 & Centred only on training, fine-tuning or architecture; at the full-text stage, also another evaluated intervention that is not prompting.\\
X5 & Secondary study, position paper or tutorial.\\
X6 & Full text with neither open nor institutional access.\\
X7 & Outside the period or the language.\\
\hline
\end{tabular}
\end{table}

\subsection{Search and records}\label{sec:busca}
The search was run on 29 October 2026, by AI, and checked by the author. The same conceptual structure was
applied to title and abstract in arXiv, IEEE Xplore and the ACM Digital Library, combining three blocks:
prompting terms, terms for the four tasks and language-model terms. The syntax and the
filters were adapted to each database (Appendix~\ref{ap:busca}). Table~\ref{tab:busca} summarizes the collection.

\begin{table}[!htb]
\caption{Records exported per database (search of 29 October 2026)}\label{tab:busca}
\centering\footnotesize
\begin{tabular}{@{}llr@{}}
\hline
Database & Collection method & Records\\
\hline
arXiv & Official API, filter by day & 1,597\\
IEEE Xplore & RIS export, filter by year & 665\\
ACM DL & Result pages, filter by month & 482\\
\hline
\multicolumn{2}{@{}l}{Total} & 2,744\\
\multicolumn{2}{@{}l}{Redundant records grouped} & 440\\
\multicolumn{2}{@{}l}{Unique studies} & 2,304\\
\hline
\end{tabular}
\end{table}

IEEE Xplore displayed 671 results; the 665 exported do not include six courses, which had no export
option. The ACM collection used the result pages because the free edition does not offer
export, and the additional complete pass of 29 October was not preserved. On 1 November 2026, the collection was
repeated and preserved, with a different ordering: the 482 DOIs matched, and the check revealed three
titles truncated by the extraction, corrected in a traceable way. The match shows the same
set on the two dates, but does not guarantee the stability of the database index.

Deduplication used DOI, arXiv identifier and normalized title, and grouped the 2,744 records into
2,304 studies, without deleting the original records. Four of these studies belong to the reference
set~\citep{zhang2011identifying}, formed by six studies from the initial analysis: S1, on
\textit{self-planning}~\citep{jiang2025selfplanningv5}; S2, on
\textit{self-debugging}~\citep{chen2023selfdebug}; S3, Reflexion~\citep{shinn2023reflexion}; S4, on
refactoring~\citep{cordeiro2024refactoring}; S5, on users of assistants~\citep{perry2023insecure};
and S6, on grounded objectives~\citep{yeo2025cogo}. S1, S2, S4 and S6 were retrieved, kept
their historical label and were left out of the random draw, without having been included or excluded. The other 2,300
are the candidates from which the calibration samples were drawn. S3 and S5 were not retrieved, because title and
abstract do not satisfy the combination of the three blocks of the search string. Since their eligibility was not assessed by the
screening, their absence indicates a potential coverage limitation, and not the confirmed loss of eligible studies.
The author decided that S1, S2, S4 and S6 would go through the batch and the full text like the others, with no automatic
entry into the corpus. S1, S4 and S6 were included; S2, only on arXiv, remained in stage 2, outside this version. S3 and S5
were left out, as a reference.

Among the 2,304 studies, 2,298 have an abstract, coming from arXiv (1,595), IEEE Xplore (446), OpenAlex
(234) and Semantic Scholar (23). Of the six without an abstract, five do not have the task in the title and were left
out of the batch; the sixth, S2160, was screened by title as a doubt, read in full text and excluded by the author (X2). Among the 2,300 candidates, 1,919 have an indication of open text, by
arXiv identifier, ACM DOI or a version located in OpenAlex, and 381 had no such indication
in the search; the access of those that went on to full text was verified at that stage.
The indication does not prove that all texts can be read, and the lack of a result in OpenAlex is not
enough to apply X6.

Criterion I4 is checked by the first public version, and not by the database date, and this check
has not yet been made for all candidates. The IEEE Xplore and ACM filters do not by themselves implement
the cut-off of 7 September 2026: 196 IEEE records have only the year 2026, and 29 ACM records have a publication
date after the cut-off, of which 15 were retrieved only in ACM in this search and depend on the search for
an earlier public version. These pending items do not exclude studies automatically.

\subsection{Screening}
Two models do the screening in separate contexts, without sharing decisions: AI-A (Claude Opus 5.5)
and AI-B (Claude Sonnet 5.5). They receive the same versioned prompt and answer include, exclude or
doubt, by title and abstract. The study goes on to full text if at least one model answers
include or doubt, so that exclusion requires both. At the full-text stage, the disagreements go to the
author, every inclusion is confirmed by the author, and only the human decision to include places a study in the
corpus. This does not constitute dual human screening.

\subsection{Screening calibration and effort measurement}\label{sec:calibracao}
Before the batch screening, the author screened two random samples of candidates, with 100 and 50 distinct
studies (seeds 20261101 and 20261102), without seeing the models' decisions. The files with those decisions
were sealed with SHA-256 before the author finished each round. The hash confirms that the files
did not change; the chronology and the absence of access depend on the record of the procedure. In the first
round, AI-A ran in the main session of the tool; in the second, in a new subagent.

A loss was defined as a study that the author included and both models excluded. If there is a loss, the
protocol requires revising the prompt and applying it to a new sample of 50 studies, with at most two revisions.
Agreement is measured by Cohen's $\kappa$ coefficient, which discounts the agreement expected by chance and
does not measure correctness. The coding was fixed before the seals were opened. In the binary $\kappa$, include and doubt count as
``goes on'', and exclude counts as ``out'', both for the author and for the union rule. The $\kappa$ between the
models uses the three categories. Table~\ref{tab:piloto} shows the results.

\begin{table}[!htb]
\caption{Results of the calibration samples}\label{tab:piloto}
\centering\footnotesize
\begin{tabular}{@{}lrr@{}}
\hline
Measure & Sample 1 & Sample 2\\
\hline
Prompt & tr-v1 & tr-v2\\
Studies that go on under the union rule & 67/100 & 38/50\\
Author's include that go on & 50/56 & 23/23\\
Author's include + doubt that go on & 67/81 & 37/41\\
Losses & 6 & 0\\
Binary $\kappa$, author $\times$ union rule & 0.645 & 0.700\\
Three-category $\kappa$, AI-A $\times$ AI-B & 0.808 & 0.613\\
\hline
\end{tabular}
\end{table}

The six losses in sample 1 triggered the revision of the prompt. Version tr-v2 started to require the answer
doubt when there is code generation by prompting and title and abstract do not settle the classification,
and made explicit that the mere presence of training or fine-tuning is not enough to apply X4. The criterion did not
change: X4 already referred to studies centred only on these topics.

Zero losses in sample 2 has a narrow meaning: none of the 23 studies that the author included was excluded
by both models. This does not prove full recall. Studies that the author marked as doubt were dropped by
exclusion by both models: eight in sample 1 and four in sample 2 (S24, S539, S962 and S2112). The author
did not record reasons in these 150 decisions, so the labels do not allow each result to be attributed
to the rules of I2 or X4. The protocol also uses the absence of loss after at most one revision as the
criterion for deciding between the full review and the fallback. This criterion was met in sample 2, which
does not validate the future process in general.

The comparison between the rounds is not a controlled experiment, for three reasons. The author already knew the
models' reasons in the second round and took part in the new rule. AI-A changed context, so as not to carry
the author's decisions. And, in the first round, the author's guide contained guidance on code generation
in broader tasks and on X4 that tr-v1 did not give to the models, so part of the losses may come
from this difference in instructions.

To calibrate the reading and estimate the effort, the author checked 12 studies in full text, with 209
pages in total, drawn at random (seed 20261103) from the 80 that went on in the two samples and have an indication
of open access. They were already in the samples, are not added to the 150 and do not enter the corpus through this
check. The author did not see the models' assessments before submitting the decisions, but the form
displayed the guidance on I2 from the tc-v1 draft, with the indication that the author could disagree with it, which may have
influenced the reading. The author included nine
studies, excluded two and was in doubt about one. The decisions coincided with those of AI-A in 10 cases and with
those of AI-B in 9, which does not measure accuracy against an external reference, and no study included by the
author was excluded by both models. The form recorded an average of 89~s per study, but the author
reported that this time does not cover the whole activity and estimated about 3~min per article. This value is a
retrospective estimate, and not a measurement.

On 2 November 2026, the two models reread the same 12 studies in new contexts, from the text extracted
from the PDFs, without figures and with compacted spaces. The recorded token consumption fell by about 40\%
in AI-A and 60\% in AI-B. AI-A kept the 12 decisions, and AI-B changed three (S2261, S427 and S306), without any
study included by the author being excluded by both models. Without repeating the reading from the PDF, it is not possible
to separate the effect of the format from the variation of the model. At that moment, the adoption of this format had not yet been
decided. On 3 November 2026, the author adopted it, with a return to the PDF in case of doubt or of a scrambled table.

In the central extrapolation, with the proportions observed in small samples, about 1,748 studies
($2{,}300 \times 38/50$) would go to full text. In the 12 texts, ten would require checking, by inclusion or
disagreement, and two were consensual exclusions, of which 20\% would be checked, which leads to about 1,515
studies. At 3~min per study, adding about 1.3~h for the 100 title and abstract exclusions, the human
checks up to the full text would require about 77~h. The projection does not cover extraction, quality, synthesis, examination
of the matrix, writing or the snowballing records, and does not discount possible unavailability among
the 381 unchecked candidates. The 12 studies are also not a simple sample of the future texts, since
they come from two samples screened with different prompts and only from studies with an indication of open access.

In the fallback variant, simulated locally on the titles and without a new search in the databases, 711 candidates remain. With the
proportion observed among the sample studies with the task in the title (37/44, with the two versions of the prompt),
about 598 would go to full text, and the same calculation gives about 27~h. However, 46 of the 79 studies included by the
author in the samples (58\%) do not have the task in the title and would be left out. This loss is the main limitation of the rapid review adopted. The ceiling of 12~h set in the protocol was removed by the author after knowing the
calibration data, which is reported as a deviation.

\subsection{Title and abstract screening}\label{sec:lote}
On 3 November 2026, AI-A and AI-B screened 715 studies by title and abstract: the 711 candidates with the task in the
title, selected locally on the titles, and S1, S2, S4 and S6. The 13 candidates from sample 2 kept
the decisions already made with tr-v2; the other 702, including the 31 from sample 1, which had tr-v1, were screened
with tr-v2 in new subagents. The decisions remained sealed with SHA-256 until the author finished the check. There was one
deviation: AI-B excluded five studies citing inclusion codes, which the response format does not accept. A new
subagent, told to cite an X code, screened them again, and one of them (S1108) went from exclude to doubt.

The protocol provided for the author to check a random sample of 100 consensual exclusions. Since there were fewer
than 100, the author checked all those that did not yet have a decision of the author's own, in a form that did not show the models'
decisions. Before the check, the author decided that a consensual exclusion marked by the author as include or
doubt would go on to full text. The rule also applies to the author's decisions in the calibration samples. This
rescue takes the study to full text, but not into the corpus.

\subsection{Full-text reading}\label{sec:textocompleto}
The studies that went on were divided into three stages: first the studies with an indication of peer review
(record in IEEE Xplore or ACM, DOI or publication reference on arXiv, or publication venue in Semantic
Scholar or OpenAlex), then the others with open access and, finally, those that depend on institutional access,
almost all from IEEE Xplore, downloaded manually by the author. Initially, the indication defined the reading order; after
the decision of 8 November 2026, it also came to delimit the studies considered in the version for the examining committee. The indication does not
prove peer review.

In stage 1, the PDFs were obtained by arXiv identifier or by DOI, and the text was extracted with Ghostscript.
For six studies, a version was found on arXiv by the exact title. ACM blocked the automated requests. In the
author's browser, 16 PDFs were downloaded before an anti-robot check, and the automation stopped at that point,
without circumventing the check. The remaining PDFs were left for manual download by the author.

Between 5 November and 7 November 2026, AI-A and AI-B read the extracted texts with tc-v2, in eight batches, returning to the PDF in case
of doubt or of a scrambled table, in new subagents. From the third batch on, each subagent received four
studies and could only read and write files, and the SHA-256 of the prompt was checked before each batch. The decisions remained sealed until the
author's check. In three studies, one of the models reported having read only part of the text: S14, without the supplementary
appendix, and S1415 and S2131, because of the context limit. In seven texts with more than 2,000 lines, it was not verified
whether the reading was complete.

The author received the studies in random order, without the models' decisions: all those with inclusion,
doubt or disagreement, and 20\% (minimum of 15) of the consensual exclusions. Because of an implementation error, the draw
applied the 20\% without the minimum, and only 3 of the 29 consensual exclusions entered. The error was noticed at import.
The author then decided to check the other 26, also without seeing the models' decisions, so that all studies
read by the models went through the author. For each study, the author recorded the decision, the criteria and a reason with
section or page. According to the author, reading the first 230 took about 12~h~30~min, done largely outside
the form. For this reason, the form's timer, which mostly measures the recording, does not serve as reading time.

During the check, the author delimited X4 for studies whose main intervention is not prompting, but which
also evaluate prompting. On 8 November 2026, the author adopted the following formulation, proposed by Claude Code from the
reasons the author had recorded: the study is included when it analyses, as a finding of its own, at least one code
result obtained by prompting in a model without fine-tuning; it is excluded by X4 when that result serves only as a
reference to show the advantage of another method. For example, S1034 and S1359 were included because they discuss
the prompting results, and S1586 was excluded because the prompt strategy, although described, serves as a
comparison for the proposed method. Having been formulated in view of the data, this boundary is reported as a
later interpretation of the criterion, and it applied to two studies included before it, S314 and S373. These two
were later excluded by X4: the author followed a recommendation by Claude Code, based on the reasons already
recorded, with no new reading recorded, and confirmed the two exclusions on 17 November 2026. At
import, two reasons that did not correspond to the recorded decision (S522 and S1586) went back to the author, who
rewrote the reasons and kept the decisions.

Between 8 November and 11 November 2026, the models read with tc-v2, under the same conditions, the 52 remaining texts of stage 1 and
the 111 of stage 3, in three batches. The author checked blind the 127 studies of the first two batches and finished that
check on 9 November 2026. Thus, 383 studies had gone through the reading by the two models and the author's check by that date. The text of S843 remains without extraction. In the 383 studies, the author included 306, was
in doubt about 9 and excluded 68. On 17 November 2026, when adjudicating the eligibility re-examination made during the extraction
(Section~\ref{sec:extracao}), the author excluded S682 by X4 and S1280 by X3, and the corpus came to 304 studies. On the same day, the author decided eight of the
nine studies in doubt, all as exclusions. For the ninth, S1617, the author decided that the date of the event counts
as the first public version, even when the DOI of the proceedings was registered later, and included it. The corpus came
to 305 studies. Also on 17 November 2026, the author finished the blind check of the 36 studies of the last batch of
stage 3, read by the models on 11 November 2026: included 28, excluded 7 and was in doubt about S1778, which was included after
examining the recommendation by Claude Code. With this, 419 studies went through the two models and the author, and the
corpus reached 334 studies. After the extraction, when adjudicating an eligibility assessment by Claude Code, the author
excluded S1721 by X4 (the quantitative result comes from fine-tuning) and S433 by X1 (arXiv version of S2053), and kept
S1929, S1886, S261, S320 and S398. The definitive corpus has 332 studies. The date of the
first public version (I4) of the included studies was checked by script, by the publication date of the first version
on arXiv and by the earliest date recorded in Crossref. All meet the period. S1610, whose DOI was registered after the
cut-off, meets it by the date of the event, under the same rule adopted for S1617.

\subsection{Stages outside this version}
The protocol provided for deciding on 10 November 2026 whether the second stage would fit within the usage limit of the models. The author brought
the decision forward to 8 November 2026: the version for the examining committee leaves out stage 2, with 233 studies only on arXiv, in line with the
cut of preprints suggested by the advisor, and reads stage 3 after stage 1. The committee corpus is thus that of the
studies with an indication of peer review, open or by institutional access. Keeping stage 3 prevents the cut
from favouring the reviewed studies that are in open access. The inclusion of stage 2 in the final version will be defined
with the advisor. The cost of the cut is losing recent work not yet published, which is reported as a
limitation.

\subsection{Extraction and quality}\label{sec:extracao}
The extraction followed the protocol revised on 10 November 2026. Each included study receives a record that registers:
\begin{itemize}
  \item tasks;
  \item techniques, with the original name, the mechanisms assigned, the role (focal or comparator) and the existence
  of a counterfactual condition;
  \item components that are not prompting, such as fine-tuning, decoding, model selection and budget;
  \item models, benchmarks, metrics and languages;
  \item risks and mitigations;
  \item discrepancies found in a targeted search (heterogeneity).
\end{itemize}
The numerical extraction targets up to three comparisons per study, chosen without considering the effect: first the
main comparison declared by the authors; in its absence, or in case of a tie, the first eligible table, the
first pertinent metric and the first model; then, the coverage of the focal techniques and of the tasks. Each value
has a literal quotation and a locator. A focal technique without a numerical result is only classified.

The instrument was calibrated twice on the 10 studies of the extraction pilot, based on the author's check. In the
first version, the author marked 43 of 85 blocks for correction; in the second, 26 of 81. The final version, ex-v2.2,
separates the mechanisms that the technique adds to the base prompt and the type of counterfactual:
\begin{itemize}
  \item \textit{package}: the whole technique against the condition without it;
  \item \textit{ingredient} or \textit{paired contrast}: an isolated mechanism;
  \item \textit{cumulative}: one round against the previous one of the same cycle;
  \item \textit{none}: no counterfactual condition.
\end{itemize}
On 17 November 2026, the author decided that a comparator system with a comparable base condition in the same table also
has an evaluation of the package, without isolation of the ingredients. The targeted revision that applied this rule changed 139
comparators in 47 studies.

Between 14 November and 16 November 2026, the frozen protocol ex-v2.2 was applied with the extracted text and the page images:
\begin{itemize}
  \item 266 studies were extracted by Claude Opus 5.5 subagents, one per study, limited to reading and writing files;
  \item 26 were extracted by Codex CLI (model gpt-6.1-sol), adopted as a second extractor after
  reproducing 18 of 21 reference comparisons in the pilot;
  \item the 10 pilot records complete 302.
\end{itemize}
Each study had a single extractor, so the extractor is a variable to be reported, and there is no dual extraction. The
raw outputs were not altered: later corrections are kept in a log with field, original value, final
value, reason and person responsible. S1610 and S1617, decided later as to I4, were extracted in the same way on
17 November 2026, and the records reached 304. The 29 studies included in the last batch of stage 3 were also extracted
this way, on the same day, and the records reached 333. The writing of S667 had
been interrupted by a safety filter of the model; at the author's request, Codex extracted it on 17 November 2026. With the
two subsequent exclusions, the corpus has 332 records, one per study. During the
extraction, the eligibility of eight studies was called into question. Claude Code re-examined them in the primary text, and the
author kept six and excluded two.

Quality was judged in the same record, by the same extractors. The rubric was adapted from
\citet{dyba2008agile}, with LLM-specific criteria~\citep{sallou2024breaking}: objective,
procedure, comparator, metrics, threats to validity, reproducibility and contamination. The possible judgements
are meets, partial, does not meet and not reported, each with a locator and a justification. The rubric qualifies the
synthesis, but does not exclude studies or produce an overall score. The judgements are made by AI and are part of the human check
described below.
% AFIRMACAO: A124
The LLM-specific part was checked against the original of \citet{sallou2024breaking}, who discuss three
threats, closed models, data leakage and reproducibility, and recommend reporting model version, seed,
prompts, query date and variability of the outputs. The reproducibility and contamination criteria
correspond to these threats, but the rubric does not ask about the query date or the variability between
runs.
% AFIRMACAO: A125
\citet{dyba2008agile} assessed each study on 11 dichotomous criteria (yes or no), grouped into
rigour, credibility and relevance, and used only the first, whether the study reports empirical research, to include or
exclude. In the rubric of this work, the objective corresponds to criterion 2 (clarity of the aims), the procedure
brings together criteria 3, 4 and 7 (context, design and data collection), the comparator corresponds to criterion 6 (control
group) and the metrics, to criterion 8 (rigour of the analysis). Criterion 1 is equivalent to I3 of the selection. Left out were
recruitment, the relationship between researcher and participants, the clarity of the findings and the value for research or
practice (criteria 5, 9, 10 and 11), which hardly apply to benchmark studies. Threats to validity are not
a criterion of Dybå and Dingsøyr and were added, and the four-level scale, with locator and justification,
replaces the dichotomous scale.

The human check has the whole study as its unit: record, comparisons and quality. The sample has
$\max(10, \lceil 0.20 \times 332\rceil) = 67$ studies:
\begin{itemize}
  \item S1608 and S1743, included because they have illegible text or tables;
  \item 59 studies from a systematic sample with a random start and a fixed seed, over the 292 records outside the pilot,
  ordered by extractor, by the existence of a manual correction and by task;
  \item 6 studies from a complement drawn by the same method over the 31 records extracted after the draw,
  when the corpus went from 304 to 334 studies; the 61 already drawn remained fixed. S1721, drawn in the complement
  and later excluded from the corpus, was replaced by S1935, by the same method.
\end{itemize}
The sample ended up with 61 studies from Opus and 6 from Codex. The protocol provided for prioritizing the studies that support
claims, but the sample was drawn before the synthesis. In compensation, each claim in the text is checked against the
original. The proportion of corrected blocks is descriptive of the
checked sample: two studies entered by obligation, those of the pilot were left out of the frame and 56 have only the author's
overall declaration, so it does not estimate an error rate for the corpus. The
check took place after the synthesis, and the status of each claim is recorded in the project repository.

\subsection{Synthesis and development of the matrix}
The synthesis starts from the records imported into the repository. The descriptive counts are generated by script:
tasks, mechanisms, counterfactuals, direction of the comparisons, quality, risks and mitigations. Distinct studies
are counted, except when the unit is the comparison. The direction of a comparison is better, equal or worse in the sense
of the metric, and is only assigned when the study defines that sense and reports both values. It describes the
extracted comparison, but does not estimate an effect or isolate a mechanism.

The synthesis by grouping of techniques was drafted in a first version by Claude Opus 5.5 subagents, over
excerpts of the records, with the rule of citing the identifier of each comparison. It was then revised and integrated by
Claude Code. Techniques are not ranked across incompatible studies, and no average is computed across
heterogeneous benchmarks. Measured result, inference of this work and suggestion of the authors are kept
separate. The contribution of a technique is discussed by what it adds to the base prompt and by the contrasts
actually evaluated.

The taxonomy and the decision-support table are examined in eight synthetic scenarios, two per task, for clarity,
grounding in the sources and concrete guidance. Claude Code wrote the scenarios before the synthesis, and they include
situations of insufficient evidence, in which the absence of a recommendation may be the appropriate answer. After the
synthesis, at the author's request, the scenarios were adjusted to fix language, domain, verification resources and
call budget, and the expectations became hypotheses. Since the adjustment took place after the synthesis, the
choice of languages and domains may reflect knowledge of the corpus. The scenarios and table v1 were
frozen with SHA-256 before the examination. For each scenario and criterion, a
judgement, a justification and ambiguities are recorded: Claude Code drafts the examination, and the author judges. After the examination, by decision
of the author, the table was revised into a version 2 over the same syntheses, and v1 remained in the repository as the object of the
examination. The advisor will review the work at the end, and the examination of the scenarios does not measure effectiveness.

\subsection{Artificial intelligence assistance}\label{sec:ia}
The search, the deduplication, the screening and the full-text reading by the models, the retrieval and extraction of the
texts, the analysis of the calibration, the import of the decisions, the counts and the draft of this text were
carried out by Claude Code (Claude Opus 5.5), with AI-B on Claude Sonnet 5.5. Codex (AI) audited
project materials, including the guide used by the author in the screening and an outline of this text. The author decided the
criteria, screened the calibration samples, checked the search and is responsible for the content. In the check of the 12 full
texts, the author did not use AI in reading the articles or in preparing the reasons, according to a declaration
recorded on 3 November 2026. On 8 November 2026, the author also declared not having used AI in the check of the 256 studies
in full text. The later change of S314 and S373 to exclusion followed a recommendation by Claude Code, with no new
reading recorded. In the extraction, Claude Opus 5.5 subagents and Codex CLI took part, as described in Section~\ref{sec:extracao}.
Claude Code carried out the targeted revision of the comparators, the eligibility re-examination, the drawing of the check
sample and the synthesis, the latter with Claude Opus 5.5 subagents. Agreement between models does not indicate correctness, since they are from the same vendor and may err
together. Assessments by AI are not treated as human review. The log of the activities, by executor and date, is
part of the project repository.
The protocol and its changes, the exact queries, the prompts, the data, the records and the logs of decisions and of
AI use are in the public project repository
(\url{https://github.com/caiofreitas1/tcc2-prompt-engineering}). The PDFs and the texts extracted from the articles are not
part of the repository.
